\documentclass{beamer}
% === collegeBeamer.sty options 可用选项 ===
%
% --- University/Institute Options 大学/机构 ---
% [polyu]  % 香港理工大学（PolyU）
% [njupt]  % 南京邮电大学 (NJUPT)
% [bnu]    % 北京师范大学 (BNU)
% 
% --- Language Options 语言 ---
% [en]     % English, plsccompile with pdflLaTeX 英语，请使用 pdflLaTeX 编译
% [zh]     % Chinese, automatically loads xeCJK and Chinese fonts, pls compile with XeLaTeX 中文，自动加载 xeCJK 包和中文字体，请使用 XeLaTeX 编译
%
% --- Example Usage ---
\usepackage[njupt, zh]{collegebeamer}
\setbeamerzoom{1}

% meta-eata
\title{南京邮电大学 Beamer 模板}
\subtitle{Using \LaTeX\ to prepare slides}
\author{15Digital Company}
\date{\today}

% document body
\begin{document}

\maketitle

\section{介绍 Introduction}

\begin{frame}
    本模板是对 \bhref{https://www.overleaf.com/latex/templates/sintef-presentation/jhbhdffczpnx}{SINTEF Presentation} 和 \bhref{https://github.com/liu-qilong/college-beamer}{College Beamer} 模板的二次创作 \vspace{\baselineskip}

    以下内容是 \bhref{mailto:federico.zenith@sintef.no}{Federico Zenith} 撰写的一篇关于如何使用 \LaTeX\ 和 beamer 制作幻灯片的简要介绍。其所有权利均归原作者所有\vspace{\baselineskip}

    本模板基于 \bhref{https://creativecommons.org/licenses/by-nc/4.0/legalcode}{Creative Commons CC BY 4.0} 许可发布
\end{frame}

\begin{frame}{用于 SINTEF 幻灯片的 Beamer}
    \begin{itemize}
        \item 我们默认你已经会使用 \LaTeX；如果你还不会，\bhref{http://en.wikibooks.org/wiki/LaTeX/}{可以在这里学习}
        \item Beamer 是 \LaTeX\ 中最流行且功能最强大的演示文稿文档类之一
        \item Beamer 还有一份详细的
              \bhref{http://www.ctan.org/tex-archive/macros/latex/contrib/beamer/doc/beameruserguide.pdf}{用户手册}
        \item 这里我们只介绍最基本的功能，帮助你快速上手
    \end{itemize}
\end{frame}

\begin{frame}{Beamer 与 PowerPoint}
    与 PowerPoint 相比，使用 \LaTeX\ 更好，因为：
    \begin{itemize}
        \item 它不是所见即所得，而是所想即所得：\\
              你只需编写内容，排版由计算机完成
        \item 生成的是 \texttt{pdf}：不会出现字体、公式、程序版本等问题
        \item 更容易保持风格、字体、强调方式等的一致性
        \item \TeX\ 的数学公式排版是最好的：
              \begin{equation*}
                  \mathrm{i}\,\hslash\frac{\partial}{\partial t} \Psi(\mathbf{r},t) =
                  -\frac{\hslash^2}{2\,m}\nabla^2\Psi(\mathbf{r},t)
                  + V(\mathbf{r})\Psi(\mathbf{r},t)
              \end{equation*}
    \end{itemize}
\end{frame}

\section{编辑 Editing}

\begin{frame}[fragile]{选择文档类}
    在图形风格上次更新之后，适用于 Beamer 的 \texttt{sintef} 主题已被更新为一个完整的文档类。
    要开始使用 \texttt{sintefbeamer}，请以下面的导言区开始一个 \LaTeX\ 文档：
    \begin{block}{最小 Beamer 文档}
        \begin{lstlisting}[language=TeX]
\documentclass{beamer}
\begin{document}

\begin{frame}{Hello, world!}
    % contents
\end {frame}

\end{document}
            \end{lstlisting}
    \end{block}
\end{frame}

\begin{frame}[fragile]{标题页}
    要设置一个典型的标题页，你需要在导言区中调用一些命令：
    \begin{block}{标题页相关命令}
        \begin{lstlisting}[language=TeX]
\title{Sample Title}
\subtitle{Sample subtitle}
\author{First Author, Second Author}
\date{Defaults to today's}\end{lstlisting}
    \end{block}
    然后你可以使用 \verb|\maketitle| 输出标题页。

    你还可以在 \verb|\maketitle| 之前使用 \verb|\titlebackground| 命令，设置不同于默认值的背景图片。

    在 \texttt{backgrounds} 文件夹中，你可以找到许多适用于 SINTEF 演示文稿标题页的标准背景图。
\end{frame}

\begin{frame}[fragile]{编写一个简单的幻灯片}
    \framesubtitle{这真的很容易！}
    \begin{itemize}[<+->]
        \item 一个典型的幻灯片通常包含项目符号列表
        \item 这些项目可以按顺序逐步显示
    \end{itemize}
    \begin{block}{包含项目列表页面的代码}<+->
        \begin{lstlisting}[language=TeX]
\begin{frame}
    \frametitle{Writing a Simple Slide}
    \framesubtitle{It's really easy!}
    \begin{itemize}[<+->]
        \item A typical slide has bulleted lists
        \item These can be uncovered in sequence
    \end{itemize}
\end{frame}\end{lstlisting}
    \end{block}
\end{frame}

\begin{frame}[fragile]{添加图片}
    添加图片的方法与普通 \LaTeX\ 中相同：
    \begin{columns}
        \begin{column}{0.7\textwidth}
            \begin{block}{添加图片的代码}
                \begin{lstlisting}[language=TeX]
\usepackage{graphicx}
% ...
\includegraphics
[width=\textwidth]{assets/demo/img.jpg}\end{lstlisting}
            \end{block}
        \end{column}
        \begin{column}{0.3\textwidth}
            \vskip5pt\includegraphics[width=\textwidth]{assets/demo/img.jpg}
        \end{column}
    \end{columns}
\end{frame}

\begin{frame}[fragile]{分栏}
    将页面分成多栏既简单又常见；通常一边放图片，另一边放文字：
    \begin{columns}
        \begin{column}{0.6\textwidth}
            这是第一栏
        \end{column}
        \begin{column}{0.3\textwidth}
            这是第二栏
        \end{column}
    \end{columns}
    \begin{block}{分栏代码}
        \begin{lstlisting}[language=TeX]
\begin{columns}
    \begin{column}{0.6\textwidth}
        This is the first column
    \end{column}
    \begin{column}{0.3\textwidth}
        And this the second
    \end{column}
    % There could be more!
\end{columns}\end{lstlisting}
    \end{block}
\end{frame}

\begin{frame}[fragile]
    \frametitle{字体}
    \begin{itemize}
        \item 字体最重要的任务是保证可读性
        \item 有一些是好的……
              \begin{itemize}
                  \item {\textrm{仅在高分辨率投影仪上使用衬线字体}}
                  \item {\textsf{否则请使用无衬线字体（或者如果你只是更喜欢它们）}}
              \end{itemize}
        \item ……也有一些不太好的：
              \begin{itemize}
                  \item {\texttt{正常正文中绝不要使用等宽字体}}
                  \item {\frakfamily{哥特体、书法体或奇怪的字体：都应始终避免使用}}
              \end{itemize}
    \end{itemize}
\end{frame}

\begin{frame}[fragile]{外观}
    \begin{itemize}
        \item 若要更改标题横线的颜色，可使用类选项 \texttt{cyandash}（默认）、\texttt{greendash}、\texttt{magentadash}、\texttt{yellowdash} 或 \texttt{nodash}。
        \item 若要插入结束页，请使用 \verb|\backmatter|。
        \item 宽高比默认为 16:9，但你也可以通过传递类选项 \texttt{aspectratio=43} 将其改为适用于老式投影仪的 4:3；Beamer 接受的其他取值也同样可以使用。
    \end{itemize}
\end{frame}

\section{表格与图表 Tables and Charts}

\begin{frame}{分组表格}
    \centering
    \begin{NJUPTtable}{>{\centering\arraybackslash}p{0.2\linewidth}|>{\centering\arraybackslash}p{0.26\linewidth}|>{\centering\arraybackslash}X}
        \NJUPTtoprule
        \textbf{领域} & \textbf{分类} & \textbf{示例方法} \\
        \NJUPTmidrule
        \multirow{2}{=}{\centering\bfseries 场景文本} & 自顶向下 & TextBoxes、EAST \\
        \cline{2-3}
        & 自底向上 & PAN++、Core-Text \\
        \NJUPTmidrule
        \multirow{2}{=}{\centering\bfseries 可解释性} & 逻辑规则 & DeepRED、KPRN \\
        \cline{2-3}
        & 决策归因 & Grad-CAM、LIME \\
        \NJUPTbottomrule
    \end{NJUPTtable}
\end{frame}

\begin{frame}{折线图示例}
    \centering
    \begin{tikzpicture}
        \begin{axis}[njupt chart, height=0.4\paperheight,
            xlabel={阶段}, ylabel={指标}, ymin=0, ymax=100, xtick={1,2,3,4,5}]
            \addplot+[thick] coordinates {(1,35) (2,48) (3,60) (4,72) (5,84)};
        \end{axis}
    \end{tikzpicture}

    \small 该示例展示单项指标在五个阶段的变化：指标从 35 稳步上升至 84。
\end{frame}

\begin{frame}{多项折线图示例}
    \centering
    \begin{tikzpicture}
        \begin{axis}[njupt chart, height=0.4\paperheight,
            xlabel={阶段}, ylabel={指标}, ymin=0, ymax=100,
            xtick={1,2,3,4,5}, legend style={at={(0.5,1.03)},anchor=south,legend columns=-1}]
            \addplot+[thick] coordinates {(1,35) (2,48) (3,60) (4,72) (5,84)};
            \addlegendentry{方案 A}
            \addplot+[thick] coordinates {(1,28) (2,39) (3,53) (4,62) (5,76)};
            \addlegendentry{方案 B}
            \addplot+[thick] coordinates {(1,22) (2,32) (3,41) (4,55) (5,66)};
            \addlegendentry{方案 C}
        \end{axis}
    \end{tikzpicture}

    \small 三条曲线便于比较不同方案的变化趋势；示例中方案 A 在各阶段均领先。
\end{frame}

\begin{frame}{趋势与构成}
    \begin{columns}[T,onlytextwidth]
        \begin{column}{0.54\textwidth}
            \centering
            \begin{tikzpicture}
                \begin{axis}[njupt chart, width=\linewidth, height=0.34\paperheight,
                    xlabel={阶段}, ylabel={指标}, ymin=0, ymax=100, xtick={1,2,3,4,5}]
                    \addplot+[thick] coordinates {(1,35) (2,48) (3,60) (4,72) (5,84)};
                \end{axis}
            \end{tikzpicture}

            \small 左图：指标随阶段稳步提高。
        \end{column}
        \begin{column}{0.42\textwidth}
            \centering
            \begin{tikzpicture}[scale=0.95]
                \fill[maincolor] (0,0) -- (90:1.5) arc[start angle=90,end angle=-90,radius=1.5] -- cycle;
                \fill[airforceblue] (0,0) -- (-90:1.5) arc[start angle=-90,end angle=-198,radius=1.5] -- cycle;
                \fill[orange!80!black] (0,0) -- (-198:1.5) arc[start angle=-198,end angle=-270,radius=1.5] -- cycle;
                \node[white,font=\small] at (0:0.85) {50\%};
                \node[white,font=\small] at (-144:0.9) {30\%};
                \node[white,font=\small] at (-234:0.9) {20\%};
            \end{tikzpicture}

            \small 右图：三个类别分别占 50\%、30\% 和 20\%。
        \end{column}
    \end{columns}
\end{frame}

\begin{frame}{柱状图示例}
    \centering
    \begin{tikzpicture}
        \begin{axis}[njupt chart, height=0.4\paperheight, ybar, bar width=16pt, ymin=0, ymax=100,
            ylabel={指标}, symbolic x coords={甲,乙,丙,丁}, xtick=data,
            nodes near coords, every node near coord/.append style={font=\scriptsize}]
            \addplot[draw=maincolor,fill=maincolor!75] coordinates {(甲,42) (乙,68) (丙,55) (丁,83)};
        \end{axis}
    \end{tikzpicture}

    \small 该示例比较四组数据，其中丁的数值最高。
\end{frame}

\section{总结 Summary}

\begin{frame}
    \frametitle{祝你好运！}
    \begin{itemize}
        \item 作为入门介绍，这些已经足够了！到现在你应该已经掌握了基本内容
        \item 如果你有修改建议或意见，\bhref{mailto:xlying23333@gmail.com}{欢迎发邮件给我！}
    \end{itemize}
\end{frame}

\begin{frame}
    \begin{block}{block}
        block
    \end{block}
    \begin{exampleblock}{example}
        example
    \end{exampleblock}
    \begin{alertblock}{alert}
        alert
    \end{alertblock}
\end{frame}

\QApage{}

\end{document}